Each new line crosses the convex sheet in a segment. If it meets the earlier lines in $k$ distinct interior points, it adds $k+1$ regions. Thus three lines create at least $1+1+1+1=4$ and at most $1+1+2+3=7$ regions. All four integer values are attainable:
\begin{itemize}
\item Four: use three parallel lines.
\item Five: let the first two cross inside the sheet, and place the third near a corner so that both its intersections with the earlier infinite lines lie outside.
\item Six: use three concurrent lines with different directions.
\item Seven: use three lines whose three pairwise intersections are distinct and inside the sheet.
\end{itemize}
The complete set is $\boxed{\{4,5,6,7\}}$.
\begin{asy}
texpreamble("\usepackage{amsmath,amssymb}");
size(240);
defaultpen(fontsize(9pt));
draw(box((-2,-2),(2,2)),gray);
draw((-2,-1)--(2,-1),blue);
draw((-2,0)--(2,0),blue);
draw((-2,1)--(2,1),blue);
label("4 regions",(0,-2.4));
\end{asy}
\begin{asy}
texpreamble("\usepackage{amsmath,amssymb}");
size(240);
defaultpen(fontsize(9pt));
draw(box((-2,-2),(2,2)),gray);
draw((0,-2)--(0,2),blue);
draw((-2,0)--(2,0),blue);
draw((1,2)--(2,1),blue);
label("5 regions",(0,-2.4));
\end{asy}
\begin{asy}
texpreamble("\usepackage{amsmath,amssymb}");
size(240);
defaultpen(fontsize(9pt));
draw(box((-2,-2),(2,2)),gray);
draw((0,-2)--(0,2),blue);
draw((-2,0)--(2,0),blue);
draw((-2,-2)--(2,2),blue);
label("6 regions",(0,-2.4));
\end{asy}
\begin{asy}
texpreamble("\usepackage{amsmath,amssymb}");
size(240);
defaultpen(fontsize(9pt));
draw(box((-2,-2),(2,2)),gray);
draw((0,-2)--(0,2),blue);
draw((-2,0)--(2,0),blue);
draw((-1,2)--(2,-1),blue);
label("7 regions",(0,-2.4));
\end{asy}
